\documentclass[../main.tex]{subfiles}

\begin{document}

\section{A Canonical Indirect/Shooting Benchmark: the Rocket Sled}\label{sec:appendix_sled}

Before tackling the ascent BVP, I validated the single-shooting code on a classical minimum-energy problem whose solution is known in closed form, the rocket sled~\cite[Lecture~5]{dclvnotes}. A unit point mass moves on a frictionless rail under a force $u$,
\begin{equation}
    \dot{r} = v, \qquad \dot{v} = u,
\end{equation}
and must be driven from rest at the origin to rest at unit displacement,
\begin{equation}
    r(0) = 0,\quad v(0) = 0, \qquad r(2) = \tfrac{1}{2},\quad v(2) = 0,
\end{equation}
in fixed time $t_f = 2$, minimizing the control energy $J = \int_0^{2} u^2\,\mathrm{d}t$. With the Hamiltonian
\begin{equation}
    \mathcal{H} = -u^2 + \lambda_r v + \lambda_v u
\end{equation}
(written in the same maximization convention used throughout this report), the costate equations are $\dot{\lambda}_r = 0$ and $\dot{\lambda}_v = -\lambda_r$, so $\lambda_r$ is constant and $\lambda_v(t) = \lambda_{v,0} - \lambda_r t$ is linear in time. Stationarity $\partial\mathcal{H}/\partial u = -2u + \lambda_v = 0$ gives the optimal control $u = \lambda_v/2$, also linear in time. Imposing the four boundary conditions fixes the two unknown initial costates analytically,
\begin{equation}
    \lambda_{r,0} = \lambda_{v,0} = \tfrac{3}{2}
    \qquad\Longrightarrow\qquad
    u^*(t) = \tfrac{3}{4}\,(1 - t),
\end{equation}
which yields $v(t) = \tfrac{3}{4}\bigl(t - t^2/2\bigr)$ and $r(t) = \tfrac{3}{4}\bigl(t^2/2 - t^3/6\bigr)$, satisfying $r(2) = \tfrac{1}{2}$ and $v(2) = 0$ exactly. Running the same \texttt{ode45}\,+\,\texttt{fsolve} shooting machinery on this two-unknown problem (\texttt{validate\_rocket\_sled.m}) recovers $(\lambda_{r,0}, \lambda_{v,0}) = (\tfrac{3}{2}, \tfrac{3}{2})$ to within $10^{-8}$ and drives the terminal residual to solver tolerance. This checks the implementation before it is applied to the ascent BVP. The rocket sled also shows how fragile shooting can be: its convergence map has sizeable regions of the initial-costate plane from which Newton iterations diverge~\cite[Lecture~5]{dclvnotes}, which is what motivates the continuation strategy of Section~\ref{sec:shooting}.

\end{document}
